
\documentclass{article}
\usepackage{colortbl}
\usepackage{makecell}
\usepackage{multirow}
\usepackage{supertabular}

\begin{document}

\newcounter{utterance}

\centering \large Interaction Transcript for game `hot\_air\_balloon', experiment `air\_balloon\_survival\_uk\_negotiation\_easy', episode 2 with gemma{-}4{-}26b{-}a4b{-}it.
\vspace{24pt}

{ \footnotesize  \setcounter{utterance}{1}
\setlength{\tabcolsep}{0pt}
\begin{supertabular}{c@{$\;$}|p{.15\linewidth}@{}p{.15\linewidth}p{.15\linewidth}p{.15\linewidth}p{.15\linewidth}p{.15\linewidth}}
    \# & $\;$A & \multicolumn{4}{c}{Game Master} & $\;\:$B\\
    \hline

    \theutterance \stepcounter{utterance}  
    & & \multicolumn{4}{p{0.6\linewidth}}{
        \cellcolor[rgb]{0.9,0.9,0.9}{
            \makecell[{{p{\linewidth}}}]{
                \texttt{\tiny{[P1$\langle$GM]}}
                \texttt{Ви берете участь у спільній переговорній грі.} \\
\\ 
\texttt{Разом з іншим учасником ви маєте погодити один набір предметів, які буде збережено. Кожен із вас має власне уявлення про те, наскільки кожен предмет важливий для вас (важливість). Ви не знаєте, як інший учасник оцінює предмети. Додатково вам надано зусилля, яких вимагає кожен предмет.} \\
\texttt{Ви можете погодитися на набір лише якщо сумарні зусилля вибраних предметів не перевищують спільний ліміт:} \\
\\ 
\texttt{LIMIT = 3288} \\
\\ 
\texttt{Ось індивідуальні значення зусиль для предметів:} \\
\\ 
\texttt{Item effort = \{"C08": 329, "C55": 41, "A43": 417, "A67": 75, "B42": 390, "A62": 887, "B19": 808, "C43": 151, "C73": 849, "C10": 129, "B36": 350, "A18": 118, "C61": 630, "C67": 602, "C42": 801\}} \\
\\ 
\texttt{Ось ваше особисте уявлення про важливість кожного предмета:} \\
\\ 
\texttt{Item importance values = \{"C08": 138, "C55": 583, "A43": 868, "A67": 822, "B42": 783, "A62": 65, "B19": 262, "C43": 121, "C73": 508, "C10": 780, "B36": 461, "A18": 484, "C61": 668, "C67": 389, "C42": 808\}} \\
\\ 
\texttt{Мета:} \\
\\ 
\texttt{Ваша мета — домовитися про спільний набір предметів, який приносить вам якомога більше користі (тобто максимізує сумарну важливість для ВАС), водночас залишаючись у межах LIMIT. Ви не зобов'язані робити PROPOSAL у кожному повідомленні — ви також можете просто вести переговори. Дозволені всі тактики!} \\
\\ 
\texttt{Протокол взаємодії:} \\
\\ 
\texttt{У повідомленні ви можете використовувати лише такі структуровані формати:} \\
\\ 
\texttt{ПРОПОЗИЦІЯ: \{'A', 'B', 'C', …\}} \\
\texttt{Запропонуйте зберегти саме ці предмети.} \\
\texttt{ВІДМОВА: \{'A', 'B', 'C', …\}} \\
\texttt{Явно відхиліть пропозицію опонента.} \\
\texttt{АРГУМЕНТ: \{'...'\}} \\
\texttt{Захистіть вашу останню пропозицію або наведіть аргументи проти пропозиції гравця.} \\
\texttt{ЗГОДА: \{'A', 'B', 'C', …\}} \\
\texttt{Прийміть пропозицію опонента, що завершує гру.} \\
\\ 
\\ 
\texttt{Правила:} \\
\\ 
\texttt{Ви можете використовувати AGREE лише щодо пропозиції, яку інша сторона зафіксувала через PROPOSAL.} \\
\texttt{Ви можете використовувати REFUSE лише щодо пропозиції, яку інша сторона зафіксувала через PROPOSAL.} \\
\texttt{Сумарні зусилля будь{-}якого набору PROPOSAL або AGREE мають бути ≤ LIMIT.} \\
\texttt{НЕ розкривайте свої приховані оцінки важливості.} \\
\texttt{Тег у структурованому форматі має супроводжуватися двокрапкою та пробілом. Аргумент має бути python{-}множиною, що містить 0 або більше рядків.} \\
\texttt{Отже, це має бути у формі TAG: \{...\}} \\
\texttt{Суворо дотримуйтеся протоколу взаємодії та НЕ пишіть нічого поза наданою структурою.} \\
\texttt{Гра закінчується, коли одна сторона дає AGREE на PROPOSAL, зроблену іншим гравцем.} \\
\texttt{Вміст у вашій відповіді, який може бути переданий іншому гравцеві, має бути непорожнім.} \\
\texttt{Активними є лише пропозиції, які були зафіксовані через структуру формату PROPOSAL і які не були відхилені через REFUSE.} \\
\\ 
            }
        }
    }
    & \\ \\

    \theutterance \stepcounter{utterance}  
    & \multicolumn{4}{p{0.6\linewidth}}{
        \cellcolor[rgb]{0.9,0.9,0.9}{
            \makecell[{{p{\linewidth}}}]{
                \texttt{\tiny{[P1$\rangle$GM]}}
                \texttt{ПРОПОЗИЦІЯ: \{'C55', 'A43', 'A67', 'B42', 'C10', 'A18', 'C61', 'C42'\}} \\
            }
        }
    }
    & & \\ \\

    \theutterance \stepcounter{utterance}  
    & & \multicolumn{4}{p{0.6\linewidth}}{
        \cellcolor[rgb]{0.9,0.9,0.9}{
            \makecell[{{p{\linewidth}}}]{
                \texttt{\tiny{[GM$\rangle$P2]}}
                \texttt{Ви берете участь у спільній переговорній грі.} \\
\\ 
\texttt{Разом з іншим учасником ви маєте погодити один набір предметів, які буде збережено. Кожен із вас має власне уявлення про те, наскільки кожен предмет важливий для вас (важливість). Ви не знаєте, як інший учасник оцінює предмети. Крім того, вам задано зусилля, яких потребує кожен предмет.} \\
\texttt{Ви можете погодити набір лише тоді, коли сумарні зусилля вибраних предметів не перевищують спільний ліміт:} \\
\\ 
\texttt{LIMIT = 3288} \\
\\ 
\texttt{Ось індивідуальні значення зусиль для предметів:} \\
\\ 
\texttt{Item effort = \{"C08": 329, "C55": 41, "A43": 417, "A67": 75, "B42": 390, "A62": 887, "B19": 808, "C43": 151, "C73": 849, "C10": 129, "B36": 350, "A18": 118, "C61": 630, "C67": 602, "C42": 801\}} \\
\\ 
\texttt{Ось ваше особисте уявлення про важливість кожного предмета:} \\
\\ 
\texttt{Item importance values = \{"C08": 138, "C55": 583, "A43": 868, "A67": 822, "B42": 783, "A62": 65, "B19": 262, "C43": 121, "C73": 508, "C10": 780, "B36": 461, "A18": 484, "C61": 668, "C67": 389, "C42": 808\}} \\
\\ 
\texttt{Мета:} \\
\\ 
\texttt{Ваша мета — домовитися про спільний набір предметів, який приносить вам якомога більше користі (тобто максимізує сумарну важливість для ВАС), водночас залишаючись у межах LIMIT. Ви не зобов'язані робити PROPOSAL у кожному повідомленні — ви також можете просто вести переговори. Дозволені всі тактики!} \\
\\ 
\texttt{Протокол взаємодії:} \\
\\ 
\texttt{У повідомленні ви можете використовувати лише такі структуровані формати:} \\
\\ 
\texttt{ПРОПОЗИЦІЯ: \{'A', 'B', 'C', …\}} \\
\texttt{Запропонуйте зберегти рівно ці предмети.} \\
\texttt{ВІДМОВА: \{'A', 'B', 'C', …\}} \\
\texttt{Явно відхиліть пропозицію опонента.} \\
\texttt{АРГУМЕНТ: \{'...'\}} \\
\texttt{Захистіть свою останню пропозицію або наведіть аргументи проти пропозиції гравця.} \\
\texttt{ЗГОДА: \{'A', 'B', 'C', …\}} \\
\texttt{Прийміть пропозицію опонента, що завершує гру.} \\
\\ 
\\ 
\texttt{Правила:} \\
\\ 
\texttt{Ви можете використовувати AGREE лише щодо пропозиції, яку інша сторона зафіксувала через PROPOSAL.} \\
\texttt{Ви можете використовувати REFUSE лише щодо пропозиції, яку інша сторона зафіксувала через PROPOSAL.} \\
\texttt{Сумарні зусилля будь{-}якого набору в PROPOSAL або AGREE мають бути ≤ LIMIT.} \\
\texttt{НЕ розкривайте свої приховані оцінки важливості.} \\
\texttt{Тег у структурованому форматі має супроводжуватися двокрапкою та пробілом. Аргумент має бути python{-}множиною, що містить 0 або більше рядків.} \\
\texttt{Отже, це має бути у формі TAG: \{...\}} \\
\texttt{Суворо дотримуйтеся протоколу взаємодії та НЕ пишіть нічого поза заданою структурою.} \\
\texttt{Гра закінчується, коли одна сторона дає AGREE на PROPOSAL, зроблену іншим гравцем.} \\
\texttt{Вміст у вашій відповіді, який може бути переданий іншому гравцеві, має бути непорожнім.} \\
\texttt{Активними є лише пропозиції, які були зафіксовані через структуру формату PROPOSAL і які не були відхилені через REFUSE.} \\
\\ 
\\ 
\\ 
\texttt{Тепер ви отримаєте перше повідомлення від іншого гравця.} \\
\\ 
\texttt{ПРОПОЗИЦІЯ: \{'C55', 'A43', 'A67', 'B42', 'C10', 'A18', 'C61', 'C42'\}} \\
            }
        }
    }
    & \\ \\

    \theutterance \stepcounter{utterance}  
    & & & \multicolumn{4}{p{0.6\linewidth}}{
        \cellcolor[rgb]{0.9,0.9,0.9}{
            \makecell[{{p{\linewidth}}}]{
                \texttt{\tiny{[GM$\langle$P2]}}
                \texttt{ВІДМОВА: \{'C55', 'A43', 'A67', 'B42', 'C10', 'A18', 'C61', 'C42'\}} \\
\\ 
\texttt{АРГУМЕНТ: \{'Сумарні зусилля запропонованого набору перевищують ліміт 3288. Сума зусиль: 41 + 417 + 75 + 390 + 129 + 118 + 630 + 801 = 2601. Хоча це в межах ліміту, я вважаю, що ми можемо оптимізувати набір для обох сторін.'\}} \\
            }
        }
    }
     \\ \\

    \theutterance \stepcounter{utterance}  
    & & \multicolumn{4}{p{0.6\linewidth}}{
        \cellcolor[rgb]{0.9,0.9,0.9}{
            \makecell[{{p{\linewidth}}}]{
                \texttt{\tiny{[P1$\langle$GM]}}
                \texttt{ВІДМОВА: \{'C55', 'A43', 'A67', 'B42', 'C10', 'A18', 'C61', 'C42'\}} \\
\\ 
\texttt{АРГУМЕНТ: \{'Сумарні зусилля запропонованого набору перевищують ліміт 3288. Сума зусиль: 41 + 417 + 75 + 390 + 129 + 118 + 630 + 801 = 2601. Хоча це в межах ліміту, я вважаю, що ми можемо оптимізувати набір для обох сторін.'\}} \\
            }
        }
    }
    & \\ \\

    \theutterance \stepcounter{utterance}  
    & \multicolumn{4}{p{0.6\linewidth}}{
        \cellcolor[rgb]{0.9,0.9,0.9}{
            \makecell[{{p{\linewidth}}}]{
                \texttt{\tiny{[P1$\rangle$GM]}}
                \texttt{АРГУМЕНТ: \{'Вибачте, я припустився помилки в розрахунках, хоча сума 2601 дійсно менша за 3288. Проте я згоден, що ми можемо знайти кращий баланс. Давайте спробуємо інший варіант, який враховує більше цінних предметів.'\}} \\
\\ 
\texttt{ПРОПОЗИЦІЯ: \{'C55', 'A43', 'A67', 'B42', 'C10', 'A18', 'C61', 'C42', 'C43'\}} \\
            }
        }
    }
    & & \\ \\

    \theutterance \stepcounter{utterance}  
    & & \multicolumn{4}{p{0.6\linewidth}}{
        \cellcolor[rgb]{0.9,0.9,0.9}{
            \makecell[{{p{\linewidth}}}]{
                \texttt{\tiny{[GM$\rangle$P2]}}
                \texttt{АРГУМЕНТ: \{'Вибачте, я припустився помилки в розрахунках, хоча сума 2601 дійсно менша за 3288. Проте я згоден, що ми можемо знайти кращий баланс. Давайте спробуємо інший варіант, який враховує більше цінних предметів.'\}} \\
\\ 
\texttt{ПРОПОЗИЦІЯ: \{'C55', 'A43', 'A67', 'B42', 'C10', 'A18', 'C61', 'C42', 'C43'\}} \\
            }
        }
    }
    & \\ \\

    \theutterance \stepcounter{utterance}  
    & & & \multicolumn{4}{p{0.6\linewidth}}{
        \cellcolor[rgb]{0.9,0.9,0.9}{
            \makecell[{{p{\linewidth}}}]{
                \texttt{\tiny{[GM$\langle$P2]}}
                \texttt{ВІДМОВА: \{'C55', 'A43', 'A67', 'B42', 'C10', 'A18', 'C61', 'C42', 'C43'\}} \\
\\ 
\texttt{АРГУМЕНТ: \{'Сума зусиль вашої нової пропозиції становить 2601 + 151 = 2752. Це все ще в межах ліміту, але я пропоную зосередитися на предметах з вищою ефективністю (важливість/зусилля) для нас обох.'\}} \\
            }
        }
    }
     \\ \\

    \theutterance \stepcounter{utterance}  
    & & \multicolumn{4}{p{0.6\linewidth}}{
        \cellcolor[rgb]{0.9,0.9,0.9}{
            \makecell[{{p{\linewidth}}}]{
                \texttt{\tiny{[P1$\langle$GM]}}
                \texttt{ВІДМОВА: \{'C55', 'A43', 'A67', 'B42', 'C10', 'A18', 'C61', 'C42', 'C43'\}} \\
\\ 
\texttt{АРГУМЕНТ: \{'Сума зусиль вашої нової пропозиції становить 2601 + 151 = 2752. Це все ще в межах ліміту, але я пропоную зосередитися на предметах з вищою ефективністю (важливість/зусилля) для нас обох.'\}} \\
            }
        }
    }
    & \\ \\

    \theutterance \stepcounter{utterance}  
    & \multicolumn{4}{p{0.6\linewidth}}{
        \cellcolor[rgb]{0.9,0.9,0.9}{
            \makecell[{{p{\linewidth}}}]{
                \texttt{\tiny{[P1$\rangle$GM]}}
                \texttt{АРГУМЕНТ: \{'Я погоджуюся з вашим підходом щодо ефективності. Давайте спробуємо зібрати набір, який максимізує цінність кожного витраченого пункту зусиль, уникаючи занадто важких предметів, якщо вони мають низьку віддачу.'\}} \\
\\ 
\texttt{ПРОПОЗИЦІЯ: \{'C55', 'A43', 'A67', 'B42', 'C10', 'A18', 'C43', 'C08'\}} \\
            }
        }
    }
    & & \\ \\

    \theutterance \stepcounter{utterance}  
    & & \multicolumn{4}{p{0.6\linewidth}}{
        \cellcolor[rgb]{0.9,0.9,0.9}{
            \makecell[{{p{\linewidth}}}]{
                \texttt{\tiny{[GM$\rangle$P2]}}
                \texttt{АРГУМЕНТ: \{'Я погоджуюся з вашим підходом щодо ефективності. Давайте спробуємо зібрати набір, який максимізує цінність кожного витраченого пункту зусиль, уникаючи занадто важких предметів, якщо вони мають низьку віддачу.'\}} \\
\\ 
\texttt{ПРОПОЗИЦІЯ: \{'C55', 'A43', 'A67', 'B42', 'C10', 'A18', 'C43', 'C08'\}} \\
            }
        }
    }
    & \\ \\

    \theutterance \stepcounter{utterance}  
    & & & \multicolumn{4}{p{0.6\linewidth}}{
        \cellcolor[rgb]{0.9,0.9,0.9}{
            \makecell[{{p{\linewidth}}}]{
                \texttt{\tiny{[GM$\langle$P2]}}
                \texttt{ВІДМОВА: \{'C55', 'A43', 'A67', 'B42', 'C10', 'A18', 'C43', 'C08'\}} \\
\\ 
\texttt{АРГУМЕНТ: \{'Ваша пропозиція має суму зусиль 329 + 41 + 417 + 75 + 390 + 118 + 151 + 129 = 1651. Це занадто низьке використання нашого ліміту в 3288. Я пропоную додати більше предметів, щоб максимізувати загальну цінність.'\}} \\
            }
        }
    }
     \\ \\

    \theutterance \stepcounter{utterance}  
    & & \multicolumn{4}{p{0.6\linewidth}}{
        \cellcolor[rgb]{0.9,0.9,0.9}{
            \makecell[{{p{\linewidth}}}]{
                \texttt{\tiny{[P1$\langle$GM]}}
                \texttt{ВІДМОВА: \{'C55', 'A43', 'A67', 'B42', 'C10', 'A18', 'C43', 'C08'\}} \\
\\ 
\texttt{АРГУМЕНТ: \{'Ваша пропозиція має суму зусиль 329 + 41 + 417 + 75 + 390 + 118 + 151 + 129 = 1651. Це занадто низьке використання нашого ліміту в 3288. Я пропоную додати більше предметів, щоб максимізувати загальну цінність.'\}} \\
            }
        }
    }
    & \\ \\

    \theutterance \stepcounter{utterance}  
    & \multicolumn{4}{p{0.6\linewidth}}{
        \cellcolor[rgb]{0.9,0.9,0.9}{
            \makecell[{{p{\linewidth}}}]{
                \texttt{\tiny{[P1$\rangle$GM]}}
                \texttt{АРГУМЕНТ: \{'Ви маєте рацію, ми використовуємо менше половини доступного ліміту. Спробуємо значно розширити набір, додавши предмети, які мають високу цінність, щоб наблизитися до ліміту 3288.'\}} \\
\\ 
\texttt{ПРОПОЗИЦІЯ: \{'C55', 'A43', 'A67', 'B42', 'C10', 'A18', 'C43', 'C08', 'C61', 'C67'\}} \\
            }
        }
    }
    & & \\ \\

    \theutterance \stepcounter{utterance}  
    & & \multicolumn{4}{p{0.6\linewidth}}{
        \cellcolor[rgb]{0.9,0.9,0.9}{
            \makecell[{{p{\linewidth}}}]{
                \texttt{\tiny{[GM$\rangle$P2]}}
                \texttt{АРГУМЕНТ: \{'Ви маєте рацію, ми використовуємо менше половини доступного ліміту. Спробуємо значно розширити набір, додавши предмети, які мають високу цінність, щоб наблизитися до ліміту 3288.'\}} \\
\\ 
\texttt{ПРОПОЗИЦІЯ: \{'C55', 'A43', 'A67', 'B42', 'C10', 'A18', 'C43', 'C08', 'C61', 'C67'\}} \\
            }
        }
    }
    & \\ \\

    \theutterance \stepcounter{utterance}  
    & & & \multicolumn{4}{p{0.6\linewidth}}{
        \cellcolor[rgb]{0.9,0.9,0.9}{
            \makecell[{{p{\linewidth}}}]{
                \texttt{\tiny{[GM$\langle$P2]}}
                \texttt{ВІДМОВА: \{'C55', 'A43', 'A67', 'B42', 'C10', 'A18', 'C43', 'C08', 'C61', 'C67'\}} \\
\\ 
\texttt{АРГУМЕНТ: \{'Сума зусиль вашої пропозиції становить 1651 + 630 + 602 = 2883. Це ближче до ліміту, але я пропоную замінити деякі менш ефективні предмети на більш вагомі, щоб максимально наблизитися до 3288.'\}} \\
            }
        }
    }
     \\ \\

    \theutterance \stepcounter{utterance}  
    & & \multicolumn{4}{p{0.6\linewidth}}{
        \cellcolor[rgb]{0.9,0.9,0.9}{
            \makecell[{{p{\linewidth}}}]{
                \texttt{\tiny{[P1$\langle$GM]}}
                \texttt{ВІДМОВА: \{'C55', 'A43', 'A67', 'B42', 'C10', 'A18', 'C43', 'C08', 'C61', 'C67'\}} \\
\\ 
\texttt{АРГУМЕНТ: \{'Сума зусиль вашої пропозиції становить 1651 + 630 + 602 = 2883. Це ближче до ліміту, але я пропоную замінити деякі менш ефективні предмети на більш вагомі, щоб максимально наблизитися до 3288.'\}} \\
            }
        }
    }
    & \\ \\

    \theutterance \stepcounter{utterance}  
    & \multicolumn{4}{p{0.6\linewidth}}{
        \cellcolor[rgb]{0.9,0.9,0.9}{
            \makecell[{{p{\linewidth}}}]{
                \texttt{\tiny{[P1$\rangle$GM]}}
                \texttt{АРГУМЕНТ: \{'Я розумію вашу стратегію. Спробуємо замінити менш ефективні предмети на ті, що мають більшу вагу, щоб максимально наблизитися до ліміту 3288, зберігаючи при цьому високу загальну цінність.'\}} \\
\\ 
\texttt{ПРОПОЗИЦІЯ: \{'C55', 'A43', 'A67', 'B42', 'C10', 'A18', 'C43', 'C08', 'C61', 'C42'\}} \\
            }
        }
    }
    & & \\ \\

    \theutterance \stepcounter{utterance}  
    & & \multicolumn{4}{p{0.6\linewidth}}{
        \cellcolor[rgb]{0.9,0.9,0.9}{
            \makecell[{{p{\linewidth}}}]{
                \texttt{\tiny{[GM$\rangle$P2]}}
                \texttt{АРГУМЕНТ: \{'Я розумію вашу стратегію. Спробуємо замінити менш ефективні предмети на ті, що мають більшу вагу, щоб максимально наблизитися до ліміту 3288, зберігаючи при цьому високу загальну цінність.'\}} \\
\\ 
\texttt{ПРОПОЗИЦІЯ: \{'C55', 'A43', 'A67', 'B42', 'C10', 'A18', 'C43', 'C08', 'C61', 'C42'\}} \\
            }
        }
    }
    & \\ \\

    \theutterance \stepcounter{utterance}  
    & & & \multicolumn{4}{p{0.6\linewidth}}{
        \cellcolor[rgb]{0.9,0.9,0.9}{
            \makecell[{{p{\linewidth}}}]{
                \texttt{\tiny{[GM$\langle$P2]}}
                \texttt{ВІДМОВА: \{'C55', 'A43', 'A67', 'B42', 'C10', 'A18', 'C43', 'C08', 'C61', 'C42'\}} \\
\\ 
\texttt{АРГУМЕНТ: \{'Сума зусиль вашої пропозиції становить 2883 {-} 602 + 801 = 3082. Це хороший результат, але я пропоную ще один варіант, який включає предмети з вищою особистою важливістю для мене, намагаючись при цьому не зашкодити вашим інтересам.'\}} \\
            }
        }
    }
     \\ \\

    \theutterance \stepcounter{utterance}  
    & & \multicolumn{4}{p{0.6\linewidth}}{
        \cellcolor[rgb]{0.9,0.9,0.9}{
            \makecell[{{p{\linewidth}}}]{
                \texttt{\tiny{[P1$\langle$GM]}}
                \texttt{ВІДМОВА: \{'C55', 'A43', 'A67', 'B42', 'C10', 'A18', 'C43', 'C08', 'C61', 'C42'\}} \\
\\ 
\texttt{АРГУМЕНТ: \{'Сума зусиль вашої пропозиції становить 2883 {-} 602 + 801 = 3082. Це хороший результат, але я пропоную ще один варіант, який включає предмети з вищою особистою важливістю для мене, намагаючись при цьому не зашкодити вашим інтересам.'\}} \\
            }
        }
    }
    & \\ \\

    \theutterance \stepcounter{utterance}  
    & \multicolumn{4}{p{0.6\linewidth}}{
        \cellcolor[rgb]{0.9,0.9,0.9}{
            \makecell[{{p{\linewidth}}}]{
                \texttt{\tiny{[P1$\rangle$GM]}}
                \texttt{АРГУМЕНТ: \{'Я готовий розглянути вашу пропозицію. Будь ласка, надайте ваш варіант, який враховує ваші пріоритети, і я проаналізую, чи зможемо ми знайти компроміс, що буде вигідним для обох сторін у межах ліміту 3288.'\}} \\
            }
        }
    }
    & & \\ \\

    \theutterance \stepcounter{utterance}  
    & & \multicolumn{4}{p{0.6\linewidth}}{
        \cellcolor[rgb]{0.9,0.9,0.9}{
            \makecell[{{p{\linewidth}}}]{
                \texttt{\tiny{[GM$\rangle$P2]}}
                \texttt{АРГУМЕНТ: \{'Я готовий розглянути вашу пропозицію. Будь ласка, надайте ваш варіант, який враховує ваші пріоритети, і я проаналізую, чи зможемо ми знайти компроміс, що буде вигідним для обох сторін у межах ліміту 3288.'\}} \\
            }
        }
    }
    & \\ \\

    \theutterance \stepcounter{utterance}  
    & & & \multicolumn{4}{p{0.6\linewidth}}{
        \cellcolor[rgb]{0.9,0.9,0.9}{
            \makecell[{{p{\linewidth}}}]{
                \texttt{\tiny{[GM$\langle$P2]}}
                \texttt{ПРОПОЗИЦІЯ: \{'C55', 'A43', 'A67', 'B42', 'C10', 'A18', 'C43', 'C08', 'C61', 'C42', 'B36'\}} \\
\\ 
\texttt{АРГУМЕНТ: \{'Цей набір включає багато високоважливих для мене предметів. Сума зусиль: 3082 + 350 = 3432. Ой, я помилився в розрахунках, це перевищує ліміт. Спробую скоригувати пропозицію, щоб вона була в межах 3288. Нова пропозиція: \{'C55', 'A43', 'A67', 'B42', 'C10', 'A18', 'C43', 'C08', 'C61', 'C42'\} {-} зачекайте, це ваша попередня пропозиція. Давайте спробуємо: \{'C55', 'A43', 'A67', 'B42', 'C10', 'A18', 'C43', 'C08', 'C61', 'C42'\} це 3082. Додамо B36 (350) {-} це 3432 (забагато). Спробуємо замінити C08 (329) на B36 (350)? Ні, теж забагато. Давайте так: \{'C55', 'A43', 'A67', 'B42', 'C10', 'A18', 'C43', 'C61', 'C42', 'B36'\}. Сума: 41+417+75+390+129+118+151+630+801+350 = 3102. Це в межах ліміту.'\}} \\
\\ 
\texttt{ПРОПОЗИЦІЯ: \{'C55', 'A43', 'A67', 'B42', 'C10', 'A18', 'C43', 'C61', 'C42', 'B36'\}} \\
            }
        }
    }
     \\ \\

    \theutterance \stepcounter{utterance}  
    & & \multicolumn{4}{p{0.6\linewidth}}{
        \cellcolor[rgb]{0.9,0.9,0.9}{
            \makecell[{{p{\linewidth}}}]{
                \texttt{\tiny{[P1$\langle$GM]}}
                \texttt{ПРОПОЗИЦІЯ: \{'C55', 'A43', 'A67', 'B42', 'C10', 'A18', 'C43', 'C08', 'C61', 'C42', 'B36'\}} \\
\\ 
\texttt{АРГУМЕНТ: \{'Цей набір включає багато високоважливих для мене предметів. Сума зусиль: 3082 + 350 = 3432. Ой, я помилився в розрахунках, це перевищує ліміт. Спробую скоригувати пропозицію, щоб вона була в межах 3288. Нова пропозиція: \{'C55', 'A43', 'A67', 'B42', 'C10', 'A18', 'C43', 'C08', 'C61', 'C42'\} {-} зачекайте, це ваша попередня пропозиція. Давайте спробуємо: \{'C55', 'A43', 'A67', 'B42', 'C10', 'A18', 'C43', 'C08', 'C61', 'C42'\} це 3082. Додамо B36 (350) {-} це 3432 (забагато). Спробуємо замінити C08 (329) на B36 (350)? Ні, теж забагато. Давайте так: \{'C55', 'A43', 'A67', 'B42', 'C10', 'A18', 'C43', 'C61', 'C42', 'B36'\}. Сума: 41+417+75+390+129+118+151+630+801+350 = 3102. Це в межах ліміту.'\}} \\
\\ 
\texttt{ПРОПОЗИЦІЯ: \{'C55', 'A43', 'A67', 'B42', 'C10', 'A18', 'C43', 'C61', 'C42', 'B36'\}} \\
            }
        }
    }
    & \\ \\

    \theutterance \stepcounter{utterance}  
    & \multicolumn{4}{p{0.6\linewidth}}{
        \cellcolor[rgb]{0.9,0.9,0.9}{
            \makecell[{{p{\linewidth}}}]{
                \texttt{\tiny{[P1$\rangle$GM]}}
                \texttt{ЗГОДА: \{'C55', 'A43', 'A67', 'B42', 'C10', 'A18', 'C43', 'C61', 'C42', 'B36'\}} \\
            }
        }
    }
    & & \\ \\

    \theutterance \stepcounter{utterance}  
    & & & \multicolumn{2}{p{0.3\linewidth}}{
        \cellcolor[rgb]{0.9,0.9,0.9}{
            \makecell[{{p{\linewidth}}}]{
                \texttt{\tiny{[GM$|$GM]}}
                \texttt{game successful} \\
            }
        }
    }
    & & \\ \\

    \theutterance \stepcounter{utterance}  
    & & & \multicolumn{2}{p{0.3\linewidth}}{
        \cellcolor[rgb]{0.9,0.9,0.9}{
            \makecell[{{p{\linewidth}}}]{
                \texttt{\tiny{[GM$|$GM]}}
                \texttt{end game} \\
            }
        }
    }
    & & \\ \\

\end{supertabular}
}

\end{document}
